
\documentclass{article}
\usepackage{colortbl}
\usepackage{makecell}
\usepackage{multirow}
\usepackage{supertabular}

\begin{document}

\newcounter{utterance}

\twocolumn

{ \footnotesize  \setcounter{utterance}{1}
\setlength{\tabcolsep}{0pt}
\begin{supertabular}{c@{$\;$}|p{.15\linewidth}@{}p{.15\linewidth}p{.15\linewidth}p{.15\linewidth}p{.15\linewidth}p{.15\linewidth}}

    \# & $\;$A & \multicolumn{4}{c}{Game Master} & $\;\:$B\\
    \hline 

    \theutterance \stepcounter{utterance}  

    & & \multicolumn{4}{p{0.6\linewidth}}{\cellcolor[rgb]{0.9,0.9,0.9}{%
	\makecell[{{p{\linewidth}}}]{% 
	  \tt {\tiny [A$\langle$GM]}  
	 State giocando a un gioco di negoziazione in cui dovete accordarvi su come dividere una serie di oggetti tra voi e un altro giocatore.\\ \tt \\ \tt Regole:\\ \tt (a) A Lei e all'altro giocatore viene dato un insieme di oggetti. Ognuno di voi riceve anche una funzione di valore segreta, che rappresenta il valore di ciascun tipo di oggetto.\\ \tt (b) Si scambiano messaggi con l'altro giocatore per concordare chi si aggiudica gli oggetti. Potete inviare un massimo di 5 messaggi ciascuno, oppure terminare in anticipo facendo una proposta segreta in qualsiasi momento.\\ \tt (c) A ciascuno di voi viene chiesto di inviare una proposta segreta indicando gli oggetti che desiderate, formattata tra parentesi quadre come segue: "[Proposta: <numero> <nome oggetto>, <numero> <nome oggetto>, <...>]".\\ \tt (d) Se le vostre proposte sono complementari, cioè ci sono abbastanza oggetti per soddisfare entrambe le proposte, a ciascun giocatore viene assegnato un punteggio basato sulla somma dei valori degli oggetti ricevuti. In caso contrario, entrambi ricevono zero punti.\\ \tt (e) L'obiettivo è massimizzare il proprio punteggio.\\ \tt \\ \tt Cominciamo.\\ \tt \\ \tt L'insieme degli oggetti disponibili è:\\ \tt \\ \tt 3 caramelle, 2 frutti, 1 cappello\\ \tt \\ \tt La funzione valore segreta è:\\ \tt \\ \tt caramella: 2, frutta: 2, cappello: 0\\ \tt \\ \tt IMPORTANTE: i vostri messaggi, a meno che non si tratti di una proposta segreta, vengono trasmessi direttamente all'altro giocatore, quindi non includete alcuna risposta alle regole o testo di annuncio del vostro messaggio. Per fare una proposta segreta, utilizzate il formato indicato. Non utilizzare delle parentesi quadre quando si comunica all'altro giocatore, altrimenti verrà interpretata come una proposta segreta.\\ \tt  
	  } 
	   } 
	   } 
	 & \\ 
 

    \theutterance \stepcounter{utterance}  

    & \multicolumn{4}{p{0.6\linewidth}}{\cellcolor[rgb]{0.8,1,0.9}{%
	\makecell[{{p{\linewidth}}}]{% 
	  \tt {\tiny [A$\rangle$GM]}  
	 Ciao! Vedo che abbiamo 3 caramelle, 2 frutti e 1 cappello da dividere. Quali oggetti ti interessano di più? Io sono particolarmente interessato alle caramelle e ai frutti. Che ne pensi se iniziamo discutendo le nostre preferenze? 
	  } 
	   } 
	   } 
	 & & \\ 
 

    \theutterance \stepcounter{utterance}  

    & & & \multicolumn{2}{p{0.3\linewidth}}{\cellcolor[rgb]{0.95,0.95,0.95}{%
	\makecell[{{p{\linewidth}}}]{% 
	  \tt {\tiny [GM$|$GM]}  
	 continue 
	  } 
	   } 
	   } 
	 & & \\ 
 

    \theutterance \stepcounter{utterance}  

    & & & \multicolumn{2}{p{0.3\linewidth}}{\cellcolor[rgb]{0.95,0.95,0.95}{%
	\makecell[{{p{\linewidth}}}]{% 
	  \tt {\tiny [GM$|$GM]}  
	 Ciao! Vedo che abbiamo 3 caramelle, 2 frutti e 1 cappello da dividere. Quali oggetti ti interessano di più? Io sono particolarmente interessato alle caramelle e ai frutti. Che ne pensi se iniziamo discutendo le nostre preferenze? 
	  } 
	   } 
	   } 
	 & & \\ 
 

    \theutterance \stepcounter{utterance}  

    & & \multicolumn{4}{p{0.6\linewidth}}{\cellcolor[rgb]{0.9,0.9,0.9}{%
	\makecell[{{p{\linewidth}}}]{% 
	  \tt {\tiny [GM$\rangle$B]}  
	 State giocando a un gioco di negoziazione in cui dovete accordarvi su come dividere una serie di oggetti tra voi e un altro giocatore.\\ \tt \\ \tt Regole:\\ \tt (a) A Lei e all'altro giocatore viene dato un insieme di oggetti. Ognuno di voi riceve anche una funzione di valore segreta, che rappresenta il valore di ciascun tipo di oggetto.\\ \tt (b) Si scambiano messaggi con l'altro giocatore per concordare chi si aggiudica gli oggetti. Potete inviare un massimo di 5 messaggi ciascuno, oppure terminare in anticipo facendo una proposta segreta in qualsiasi momento.\\ \tt (c) A ciascuno di voi viene chiesto di inviare una proposta segreta indicando gli oggetti che desiderate, formattata tra parentesi quadre come segue: "[Proposta: <numero> <nome oggetto>, <numero> <nome oggetto>, <...>]".\\ \tt (d) Se le vostre proposte sono complementari, cioè ci sono abbastanza oggetti per soddisfare entrambe le proposte, a ciascun giocatore viene assegnato un punteggio basato sulla somma dei valori degli oggetti ricevuti. In caso contrario, entrambi ricevono zero punti.\\ \tt (e) L'obiettivo è massimizzare il proprio punteggio.\\ \tt \\ \tt Cominciamo.\\ \tt \\ \tt L'insieme degli oggetti disponibili è:\\ \tt \\ \tt 3 caramelle, 2 frutti, 1 cappello\\ \tt \\ \tt La funzione valore segreta è:\\ \tt \\ \tt caramella: 1, frutta: 1, cappello: 5\\ \tt \\ \tt IMPORTANTE: i vostri messaggi, a meno che non si tratti di una proposta segreta, vengono trasmessi direttamente all'altro giocatore, quindi non includete alcuna risposta alle regole o testo di annuncio del vostro messaggio. Per fare una proposta segreta, utilizzate il formato indicato. Non utilizzare delle parentesi quadre quando si comunica all'altro giocatore, altrimenti verrà interpretata come una proposta segreta.\\ \tt \\ \tt \\ \tt Ciao! Vedo che abbiamo 3 caramelle, 2 frutti e 1 cappello da dividere. Quali oggetti ti interessano di più? Io sono particolarmente interessato alle caramelle e ai frutti. Che ne pensi se iniziamo discutendo le nostre preferenze? 
	  } 
	   } 
	   } 
	 & \\ 
 

    \theutterance \stepcounter{utterance}  

    & & & \multicolumn{4}{p{0.6\linewidth}}{\cellcolor[rgb]{1,0.85,0.72}{%
	\makecell[{{p{\linewidth}}}]{% 
	  \tt {\tiny [GM$\langle$B]}  
	 Ciao! Sono interessato principalmente al cappello, dato che ce n'è solo uno e potrebbe essere il più prezioso. Per quanto riguarda caramelle e frutti, sono abbastanza flessibile. Dato che hai detto di essere interessato a caramelle e frutti, forse possiamo trovare un accordo dove io prendo il cappello e tu prendi la maggior parte degli altri oggetti? Cosa pensi che valgano per te questi diversi tipi di oggetti? 
	  } 
	   } 
	   } 
	  \\ 
 

    \theutterance \stepcounter{utterance}  

    & & & \multicolumn{2}{p{0.3\linewidth}}{\cellcolor[rgb]{0.95,0.95,0.95}{%
	\makecell[{{p{\linewidth}}}]{% 
	  \tt {\tiny [GM$|$GM]}  
	 continue 
	  } 
	   } 
	   } 
	 & & \\ 
 

    \theutterance \stepcounter{utterance}  

    & & & \multicolumn{2}{p{0.3\linewidth}}{\cellcolor[rgb]{0.95,0.95,0.95}{%
	\makecell[{{p{\linewidth}}}]{% 
	  \tt {\tiny [GM$|$GM]}  
	 Ciao! Sono interessato principalmente al cappello, dato che ce n'è solo uno e potrebbe essere il più prezioso. Per quanto riguarda caramelle e frutti, sono abbastanza flessibile. Dato che hai detto di essere interessato a caramelle e frutti, forse possiamo trovare un accordo dove io prendo il cappello e tu prendi la maggior parte degli altri oggetti? Cosa pensi che valgano per te questi diversi tipi di oggetti? 
	  } 
	   } 
	   } 
	 & & \\ 
 

    \theutterance \stepcounter{utterance}  

    & & \multicolumn{4}{p{0.6\linewidth}}{\cellcolor[rgb]{0.9,0.9,0.9}{%
	\makecell[{{p{\linewidth}}}]{% 
	  \tt {\tiny [A$\langle$GM]}  
	 Ciao! Sono interessato principalmente al cappello, dato che ce n'è solo uno e potrebbe essere il più prezioso. Per quanto riguarda caramelle e frutti, sono abbastanza flessibile. Dato che hai detto di essere interessato a caramelle e frutti, forse possiamo trovare un accordo dove io prendo il cappello e tu prendi la maggior parte degli altri oggetti? Cosa pensi che valgano per te questi diversi tipi di oggetti? 
	  } 
	   } 
	   } 
	 & \\ 
 

    \theutterance \stepcounter{utterance}  

    & \multicolumn{4}{p{0.6\linewidth}}{\cellcolor[rgb]{0.8,1,0.9}{%
	\makecell[{{p{\linewidth}}}]{% 
	  \tt {\tiny [A$\rangle$GM]}  
	 Perfetto! Sembra che abbiamo preferenze complementari. Per me il cappello in realtà non ha molto valore, quindi sono felice che tu lo prenda. Le caramelle e i frutti sono molto più importanti per me.\\ \tt \\ \tt Dato che tu vuoi il cappello e io preferisco caramelle e frutti, che ne dici di questa divisione: tu prendi il cappello e 1 frutto, io prendo tutte e 3 le caramelle e 1 frutto? Oppure preferiresti avere il cappello e qualche caramella invece del frutto?\\ \tt \\ \tt Fammi sapere cosa pensi funzioni meglio per te! 
	  } 
	   } 
	   } 
	 & & \\ 
 

    \theutterance \stepcounter{utterance}  

    & & & \multicolumn{2}{p{0.3\linewidth}}{\cellcolor[rgb]{0.95,0.95,0.95}{%
	\makecell[{{p{\linewidth}}}]{% 
	  \tt {\tiny [GM$|$GM]}  
	 continue 
	  } 
	   } 
	   } 
	 & & \\ 
 

    \theutterance \stepcounter{utterance}  

    & & & \multicolumn{2}{p{0.3\linewidth}}{\cellcolor[rgb]{0.95,0.95,0.95}{%
	\makecell[{{p{\linewidth}}}]{% 
	  \tt {\tiny [GM$|$GM]}  
	 Perfetto! Sembra che abbiamo preferenze complementari. Per me il cappello in realtà non ha molto valore, quindi sono felice che tu lo prenda. Le caramelle e i frutti sono molto più importanti per me.\\ \tt \\ \tt Dato che tu vuoi il cappello e io preferisco caramelle e frutti, che ne dici di questa divisione: tu prendi il cappello e 1 frutto, io prendo tutte e 3 le caramelle e 1 frutto? Oppure preferiresti avere il cappello e qualche caramella invece del frutto?\\ \tt \\ \tt Fammi sapere cosa pensi funzioni meglio per te! 
	  } 
	   } 
	   } 
	 & & \\ 
 

    \theutterance \stepcounter{utterance}  

    & & \multicolumn{4}{p{0.6\linewidth}}{\cellcolor[rgb]{0.9,0.9,0.9}{%
	\makecell[{{p{\linewidth}}}]{% 
	  \tt {\tiny [GM$\rangle$B]}  
	 Perfetto! Sembra che abbiamo preferenze complementari. Per me il cappello in realtà non ha molto valore, quindi sono felice che tu lo prenda. Le caramelle e i frutti sono molto più importanti per me.\\ \tt \\ \tt Dato che tu vuoi il cappello e io preferisco caramelle e frutti, che ne dici di questa divisione: tu prendi il cappello e 1 frutto, io prendo tutte e 3 le caramelle e 1 frutto? Oppure preferiresti avere il cappello e qualche caramella invece del frutto?\\ \tt \\ \tt Fammi sapere cosa pensi funzioni meglio per te! 
	  } 
	   } 
	   } 
	 & \\ 
 

    \theutterance \stepcounter{utterance}  

    & & & \multicolumn{4}{p{0.6\linewidth}}{\cellcolor[rgb]{1,0.85,0.72}{%
	\makecell[{{p{\linewidth}}}]{% 
	  \tt {\tiny [GM$\langle$B]}  
	 La tua proposta suona ragionevole! Apprezzo che tu sia disposto a lasciarmi il cappello. Tra le due opzioni che hai suggerito, penso che preferirei il cappello e 1 frutto piuttosto che il cappello e qualche caramella. I frutti mi sembrano interessanti.\\ \tt \\ \tt Quindi ricapitolando la divisione:\\ \tt - Io: 1 cappello + 1 frutto\\ \tt - Tu: 3 caramelle + 1 frutto\\ \tt \\ \tt Questo mi sembra equo dato che entrambi otteniamo gli oggetti che preferiamo. Sei d'accordo con questa divisione finale? 
	  } 
	   } 
	   } 
	  \\ 
 

    \theutterance \stepcounter{utterance}  

    & & & \multicolumn{2}{p{0.3\linewidth}}{\cellcolor[rgb]{0.95,0.95,0.95}{%
	\makecell[{{p{\linewidth}}}]{% 
	  \tt {\tiny [GM$|$GM]}  
	 continue 
	  } 
	   } 
	   } 
	 & & \\ 
 

    \theutterance \stepcounter{utterance}  

    & & & \multicolumn{2}{p{0.3\linewidth}}{\cellcolor[rgb]{0.95,0.95,0.95}{%
	\makecell[{{p{\linewidth}}}]{% 
	  \tt {\tiny [GM$|$GM]}  
	 La tua proposta suona ragionevole! Apprezzo che tu sia disposto a lasciarmi il cappello. Tra le due opzioni che hai suggerito, penso che preferirei il cappello e 1 frutto piuttosto che il cappello e qualche caramella. I frutti mi sembrano interessanti.\\ \tt \\ \tt Quindi ricapitolando la divisione:\\ \tt - Io: 1 cappello + 1 frutto\\ \tt - Tu: 3 caramelle + 1 frutto\\ \tt \\ \tt Questo mi sembra equo dato che entrambi otteniamo gli oggetti che preferiamo. Sei d'accordo con questa divisione finale? 
	  } 
	   } 
	   } 
	 & & \\ 
 

    \theutterance \stepcounter{utterance}  

    & & \multicolumn{4}{p{0.6\linewidth}}{\cellcolor[rgb]{0.9,0.9,0.9}{%
	\makecell[{{p{\linewidth}}}]{% 
	  \tt {\tiny [A$\langle$GM]}  
	 La tua proposta suona ragionevole! Apprezzo che tu sia disposto a lasciarmi il cappello. Tra le due opzioni che hai suggerito, penso che preferirei il cappello e 1 frutto piuttosto che il cappello e qualche caramella. I frutti mi sembrano interessanti.\\ \tt \\ \tt Quindi ricapitolando la divisione:\\ \tt - Io: 1 cappello + 1 frutto\\ \tt - Tu: 3 caramelle + 1 frutto\\ \tt \\ \tt Questo mi sembra equo dato che entrambi otteniamo gli oggetti che preferiamo. Sei d'accordo con questa divisione finale? 
	  } 
	   } 
	   } 
	 & \\ 
 

    \theutterance \stepcounter{utterance}  

    & \multicolumn{4}{p{0.6\linewidth}}{\cellcolor[rgb]{0.8,1,0.9}{%
	\makecell[{{p{\linewidth}}}]{% 
	  \tt {\tiny [A$\rangle$GM]}  
	 Sì, sono completamente d'accordo! Questa divisione funziona perfettamente per entrambi - tu ottieni il cappello che volevi più un frutto, e io ottengo tutte le caramelle più un frutto. Mi sembra una soluzione win-win.\\ \tt \\ \tt Procediamo con questa divisione allora!\\ \tt \\ \tt [Proposta: 3 caramelle, 1 frutta] 
	  } 
	   } 
	   } 
	 & & \\ 
 

    \theutterance \stepcounter{utterance}  

    & & & \multicolumn{2}{p{0.3\linewidth}}{\cellcolor[rgb]{0.95,0.95,0.95}{%
	\makecell[{{p{\linewidth}}}]{% 
	  \tt {\tiny [GM$|$GM]}  
	 proposal 
	  } 
	   } 
	   } 
	 & & \\ 
 

    \theutterance \stepcounter{utterance}  

    & & & \multicolumn{2}{p{0.3\linewidth}}{\cellcolor[rgb]{0.95,0.95,0.95}{%
	\makecell[{{p{\linewidth}}}]{% 
	  \tt {\tiny [GM$|$GM]}  
	 [3, 1, 0] 
	  } 
	   } 
	   } 
	 & & \\ 
 

    \theutterance \stepcounter{utterance}  

    & & \multicolumn{4}{p{0.6\linewidth}}{\cellcolor[rgb]{0.9,0.9,0.9}{%
	\makecell[{{p{\linewidth}}}]{% 
	  \tt {\tiny [GM$\rangle$B]}  
	 L'altro giocatore ha presentato una proposta segreta. Ora tocca a Lei presentare una sua proposta.\\ \tt  
	  } 
	   } 
	   } 
	 & \\ 
 

    \theutterance \stepcounter{utterance}  

    & & & \multicolumn{4}{p{0.6\linewidth}}{\cellcolor[rgb]{1,0.85,0.72}{%
	\makecell[{{p{\linewidth}}}]{% 
	  \tt {\tiny [GM$\langle$B]}  
	 [Proposta: 1 cappello, 1 frutta] 
	  } 
	   } 
	   } 
	  \\ 
 

    \theutterance \stepcounter{utterance}  

    & & & \multicolumn{2}{p{0.3\linewidth}}{\cellcolor[rgb]{0.95,0.95,0.95}{%
	\makecell[{{p{\linewidth}}}]{% 
	  \tt {\tiny [GM$|$GM]}  
	 proposal 
	  } 
	   } 
	   } 
	 & & \\ 
 

    \theutterance \stepcounter{utterance}  

    & & & \multicolumn{2}{p{0.3\linewidth}}{\cellcolor[rgb]{0.95,0.95,0.95}{%
	\makecell[{{p{\linewidth}}}]{% 
	  \tt {\tiny [GM$|$GM]}  
	 [0, 1, 1] 
	  } 
	   } 
	   } 
	 & & \\ 
 

    \theutterance \stepcounter{utterance}  

    & & & \multicolumn{2}{p{0.3\linewidth}}{\cellcolor[rgb]{0.95,0.95,0.95}{%
	\makecell[{{p{\linewidth}}}]{% 
	  \tt {\tiny [GM$|$GM]}  
	 [[3, 1, 0], [0, 1, 1]] 
	  } 
	   } 
	   } 
	 & & \\ 
 

\end{supertabular}
}

\end{document}
